\chapter{The six operations with finite coefficients}
\label{ch:six-operations}

This chapter constructs the functorial language in which the trace formula and
duality will be stated.  Throughout, \(S\) is a noetherian scheme on which the
integer \(m\) is invertible, and
\[
  \Lambda=\mathbf Z/m\mathbf Z.
\]
All schemes are separated and of finite type over \(S\), unless a different
hypothesis is displayed.  The symbol
\(D^b_c(X,\Lambda)\) denotes the bounded derived category of complexes of
etale \(\Lambda\)-modules whose cohomology sheaves are constructible.

There are two logically distinct assertions in the phrase ``six operations.''
First, the operations exist on the ambient derived categories of sheaves.
Second, under finite-type hypotheses they preserve bounded constructible
complexes.  The second assertion contains the finiteness and base-change
theorems and is the difficult one.  Keeping these assertions separate is
essential both mathematically and in a formalization.

\section{The elementary four operations}

Let \(f:X\to Y\) be a morphism.  The inverse-image functor
\[
  f^*:\operatorname{Mod}(Y_{\mathrm{et}},\Lambda)
       \longrightarrow
       \operatorname{Mod}(X_{\mathrm{et}},\Lambda)
\]
is exact.  Its right adjoint \(f_*\) is left exact, and therefore has a right
derived functor
\[
  Rf_*:D^+(X_{\mathrm{et}},\Lambda)
       \longrightarrow D^+(Y_{\mathrm{et}},\Lambda).
\]
The sheaf tensor product is right exact in each variable, and internal Hom is
left exact in its second variable.  Their derived versions are
\[
  -\otimes_\Lambda^L-,
  \qquad
  R\mathcal Hom_\Lambda(-,-).
\]
For \(K,L,M\) in the ranges where the expressions are defined, the fundamental
adjunction is
\[
  R\mathcal Hom(K\otimes_\Lambda^L L,M)
  \simeq
  R\mathcal Hom\bigl(K,R\mathcal Hom(L,M)\bigr).
\]
It is an isomorphism of functors, not an equality of chosen complexes.

\begin{proposition}[Pullback, tensor, and internal Hom]
\label{prop:tensor-hom-pullback}
For every morphism \(f:X\to Y\) there are canonical isomorphisms
\[
  Lf^*(K\otimes_\Lambda^L L)
   \simeq Lf^*K\otimes_\Lambda^L Lf^*L
\]
and, if \(K\) is locally represented by a bounded complex of finite free
\(\Lambda\)-modules,
\[
  Lf^*R\mathcal Hom_\Lambda(K,L)
   \simeq
  R\mathcal Hom_\Lambda(Lf^*K,Lf^*L).
\]
Here \(Lf^*=f^*\), since inverse image is exact on etale sheaves.
\end{proposition}

\begin{proof}
The tensor assertion follows already before deriving: inverse image is
obtained by sheafifying the objectwise inverse image, and both sheafification
and inverse image commute with finite colimits and tensor products.  Resolve
one tensor factor by flat sheaves to obtain the derived statement.  For
internal Hom, work etale locally where \(K\) is a bounded finite free complex.
The claim then reduces term by term to
\(f^*\mathcal Hom(\Lambda^r,L)=\mathcal Hom(\Lambda^r,f^*L)\), and totalize the
finite double complex.
\end{proof}

\section{Extension by zero and compactifiable morphisms}

If \(j:U\hookrightarrow X\) is an open immersion, extension by zero is the
exact functor \(j_!\) characterized by
\[
  (j_!\mathcal F)_{\bar x}=
  \begin{cases}
    \mathcal F_{\bar x},&\bar x\text{ lies over }U,\\
    0,&\bar x\text{ lies over }X\setminus U.
  \end{cases}
\]
It is left adjoint to \(j^*\).  If \(i:Z\hookrightarrow X\) is the
complementary closed immersion, then every \(K\in D(X_{\mathrm{et}},\Lambda)\)
fits into the localization triangle
\[
  j_!j^*K\longrightarrow K\longrightarrow i_*i^*K
  \longrightarrow j_!j^*K[1].
\]

A separated morphism \(f:X\to Y\) of finite type is compactifiable:
Nagata compactification gives a factorization
\[
  X\xrightarrow{j}\overline X\xrightarrow{\bar f}Y
\]
with \(j\) open and \(\bar f\) proper.

\begin{definition}[Exceptional direct image]
\label{def:Rfshriek}
For a chosen compactification \(f=\bar f j\), put
\[
  Rf_!K:=R\bar f_*j_!K.
\]
\end{definition}

The notation is legitimate only after proving independence of the chosen
compactification.  This point cannot be hidden in the definition.

\begin{theorem}[Independence of compactification]
\label{thm:Rfshriek-independence}
Let \(f:X\to Y\) be separated and of finite type.  The complexes
\(R\bar f_*j_!K\), as \((j,\bar f)\) ranges over compactifications of \(f\),
are related by unique canonical isomorphisms satisfying the cocycle condition.
Consequently \(Rf_!\) is a well-defined functor up to a contractible system of
canonical choices.  It is functorial under composition:
\[
  R(gf)_!\simeq Rg_!\,Rf_!.
\]
\end{theorem}

\begin{proof}[Proof capsule]
Given two compactifications, take the schematic closure of the diagonal image
of \(X\) in their fiber product.  After replacing it by the closure, one gets
a third compactification dominating both by proper maps \(p_1,p_2\) that are
the identity on \(X\).  For such a proper modification,
\[
  Rp_{a*}j'_!K\simeq j_{a!}K.
\]
This is checked on geometric stalks.  Over \(X\) the map is the identity; at a
point of the boundary, proper base change identifies the stalk with the
cohomology of a proper fiber carrying a sheaf whose support is empty, hence
with zero.  The comparison for a common refinement is therefore canonical.
Any two common refinements admit a third, so the comparisons obey the cocycle
condition.  For composition, compactify successively and apply the same
comparison.  This is the compactification theorem of
\cite[Expos\'e XVII]{SGA4}.
\end{proof}

\begin{proposition}[Localization for \(Rf_!\)]
\label{prop:six-localization}
Let \(i:Z\hookrightarrow X\) be closed with open complement \(j:U\hookrightarrow
X\), and let \(f:X\to Y\) be separated of finite type.  Then
\[
  R(fj)_!j^*K\longrightarrow Rf_!K
  \longrightarrow R(fi)_!i^*K\longrightarrow
\]
is a distinguished triangle.
\end{proposition}

\begin{proof}
Apply \(Rf_!\) to the localization triangle.  The composition isomorphism in
Theorem~\ref{thm:Rfshriek-independence} identifies the first and third terms
with the displayed ones.
\end{proof}

\section{Proper base change without circularity}

Consider a cartesian square
\[
\begin{tikzcd}
X' \arrow[r,"g'"] \arrow[d,"f'"'] & X \arrow[d,"f"]\\
Y' \arrow[r,"g"'] & Y .
\end{tikzcd}
\]
Adjunction gives the base-change morphism
\[
  g^*Rf_*K\longrightarrow Rf'_*g'^*K.
\]

\begin{theorem}[Proper base change, finite coefficients]
\label{thm:proper-base-change-torsion}
If \(f\) is proper and \(K\in D^b_c(X,\Lambda)\), the base-change morphism is
an isomorphism.  Equivalently, for every geometric point \(\bar y\to Y\),
\[
  (R^rf_*K)_{\bar y}
   \simeq
  H^r(X_{\bar y,\mathrm{et}},K|_{X_{\bar y}})
\]
for all \(r\).
\end{theorem}

\begin{proof}[Dependency-accurate proof capsule]
The theorem is proved simultaneously with constructible finiteness, not by
using the finished finiteness theorem as an earlier black box.  In relative
dimension at most \(d\), let \(A_d\) denote the affine cohomological bound,
\(P_d\) proper base change, and \(C_d\) constructibility of direct images.
The induction is
\[
  \{A_e,P_e,C_e:e<d\}\Longrightarrow A_d
  \Longrightarrow P_d\Longrightarrow C_d.
\]
The preliminary Artin bound of Chapter~\ref{ch:etale-foundations}, in
particular \eqref{eq:artin-global-bound}, is independent of proper base change
and supplies the boundedness needed for all spectral sequences below.

For \(P_d\), the assertion is local for the strict henselian topology on
\(Y\).  Thus take \(Y\) strictly local with closed geometric point \(\bar y\);
one must prove
\[
  R\Gamma(X,K)\longrightarrow R\Gamma(X_{\bar y},K)
\]
is an isomorphism.  After noetherian devissage on the support and alteration
of the generic stratum, construct a proper hypercover
\(X_\bullet\to X\) whose terms are assembled from schemes for which the
claim is either finite, smooth with a normal-crossings boundary, or of smaller
relative dimension.  Cohomological descent gives spectral sequences
\[
 E_1^{a,b}=H^b(X_a,K_a)\Rightarrow H^{a+b}(X,K),
\qquad
 {E'}_1^{a,b}=H^b((X_a)_{\bar y},K_a)
 \Rightarrow H^{a+b}(X_{\bar y},K).
\]
Only finitely many rows can contribute in any total degree, by the independent
Artin bound.  The comparison on the \(E_1\)-page is an isomorphism by the
lower-dimensional induction and the elementary finite and smooth local
cases.  Hence the abutment comparison is an isomorphism.  Devissage through a
finite stratification extends the conclusion from locally constant sheaves to
all constructible complexes.  Once \(P_d\) is known, the stalk formula and
constructibility on strata prove \(C_d\).  This is the organization of the
finiteness proof in \cite[Expos\'es XIV--XVI]{SGA4}; it avoids both
\(P_d\Rightarrow P_d\) and \(C_d\Rightarrow P_d\).
\end{proof}

\begin{theorem}[Smooth base change]
\label{thm:smooth-base-change}
In the cartesian square above, suppose that \(g\) is smooth.  For every
quasi-compact quasi-separated \(f\) and every torsion constructible complex
\(K\), the natural map
\[
  g^*Rf_*K\longrightarrow Rf'_*g'^*K
\]
is an isomorphism whenever the displayed direct images are defined in the
bounded constructible range.
\end{theorem}

\begin{proof}[Proof capsule]
The statement is local on \(X\) and \(Y\) and on the geometric stalk of
\(Y'\).  Smooth morphisms admit etale-local coordinate charts to affine
space.  Homotopy invariance for torsion prime to the residue characteristics,
together with the acyclicity of the strict henselian local model, proves the
claim for the projection \(\mathbf A^r_Y\to Y\).  Etale descent transfers it
to an arbitrary smooth \(g\).  A finite affine cover of \(X\) and its
Cech-to-derived spectral sequence remove the affineness assumption.  The
cohomological bounds used for convergence are the Artin bounds, not proper
base change.  See \cite[Expos\'e XVI]{SGA4}.
\end{proof}

\begin{theorem}[Base change for compact support]
\label{thm:compact-support-base-change}
For any cartesian square as above with \(f\) separated of finite type,
\[
  g^*Rf_!K\xrightarrow{\sim}Rf'_!g'^*K
\]
for \(K\in D^b_c(X,\Lambda)\).
\end{theorem}

\begin{proof}
Choose a compactification \(f=\bar f j\) and base-change it to \(Y'\).
Inverse image commutes with extension by zero, while proper base change gives
\[
g^*R\bar f_*j_!K
\simeq R\bar f'_*g_{\overline X}^*j_!K
\simeq R\bar f'_*j'_!g'^*K.
\]
Independence of compactification makes the isomorphism canonical.
\end{proof}

\section{Projection formula and Kunneth}

\begin{theorem}[Projection formula]
\label{thm:projection-formula-torsion}
Let \(f:X\to Y\) be separated of finite type.  For
\(K\in D^b_c(X,\Lambda)\) and \(L\in D^b_c(Y,\Lambda)\), with one factor of
finite Tor dimension, the canonical morphism
\[
  Rf_!K\otimes_\Lambda^L L
  \longrightarrow
  Rf_!(K\otimes_\Lambda^L f^*L)
\]
is an isomorphism.
\end{theorem}

\begin{proof}
For an open immersion this is checked on stalks.  For a proper morphism it is
the derived projection formula: resolve \(L\) by flat sheaves, use adjunction,
and reduce to the underived formula on injective resolutions.  Factor \(f\)
as open followed by proper and combine the two cases.
\end{proof}

\begin{theorem}[Kunneth formula with compact support]
\label{thm:kunneth-torsion}
Let \(X,Y\) be separated of finite type over a field \(k\), and let
\(K\in D^b_c(X,\Lambda)\), \(L\in D^b_c(Y,\Lambda)\).  With
\(p_X,p_Y\) the projections,
\[
 R\Gamma_c(X,K)\otimes_\Lambda^L R\Gamma_c(Y,L)
 \xrightarrow{\sim}
 R\Gamma_c\!\left(X\times_kY,
 p_X^*K\otimes_\Lambda^L p_Y^*L\right).
\]
\end{theorem}

\begin{proof}
Let \(a_X:X\to\operatorname{Spec}k\) and similarly for \(Y\).  Apply compact
support base change to the square defined by \(a_X\) and \(Y\), then the
projection formula:
\[
\begin{aligned}
R(a_X\times a_Y)_!
 (p_X^*K\otimes^Lp_Y^*L)
&\simeq
R(a_Y)_!\bigl(Rp_{Y!}p_X^*K\otimes^LL\bigr)\\
&\simeq
R(a_Y)_!\bigl(a_Y^*R(a_X)_!K\otimes^LL\bigr)\\
&\simeq R(a_X)_!K\otimes^LR(a_Y)_!L.
\end{aligned}
\]
The structure map to the point identifies these terms with compactly
supported cohomology.
\end{proof}

\section{The finite-coefficient package}

For later reference, the outcome can be stated as follows.

\begin{theorem}[Finite-coefficient six-operation package]
For separated finite-type schemes over \(S\), the categories
\(D^b_c(X,\Lambda)\) carry functors
\[
  Rf_*,\quad Rf_!,\quad f^*,\quad Rf^!,
  \quad -\otimes_\Lambda^L-,\quad R\mathcal Hom_\Lambda(-,-),
\]
with the standard adjunctions, composition isomorphisms, localization
triangles, base-change isomorphisms, projection formula, and Kunneth formula.
Here \(Rf^!\), right adjoint to \(Rf_!\), is constructed by Brown
representability on the ambient derived category and shown to preserve the
constructible range by duality; its geometric description is deferred to
Chapter~\ref{ch:duality}.  All coherence diagrams commute.
\end{theorem}

\begin{proof}
The first five operations and their compatibilities were constructed above.
The exceptional inverse image and its constructibility are the duality
theorem of \cite[Expos\'e XVIII]{SGA4}.  Coherence is first checked on the
underived adjunctions and on a common refinement of compactifications, then
passes to derived functors.  Since every comparison is characterized by an
adjunction or a universal property, two composites around any coherence
diagram have the same adjoint and are equal.
\end{proof}

\section{Formalization checklist}

A formal library should expose the following objects rather than one opaque
``six operations'' assumption.
\begin{enumerate}
  \item Categories \(D^b_c(X,\Lambda)\), together with the inclusion into the
  ambient derived category and a predicate for constructibility.
  \item Functors \(f^*,Rf_*,Rf_!,Rf^!,\otimes^L,R\mathcal Hom\), each with its
  own hypotheses and adjunction data.
  \item Compactifications as a cofiltered category; comparison isomorphisms
  for refinements and their cocycle law.
  \item Named natural transformations for base change and projection formula,
  followed by separate theorems asserting that they are isomorphisms.
  \item The simultaneous induction data \(A_d,P_d,C_d\).  Its dependency graph
  should be represented explicitly so that constructible finiteness cannot be
  imported into the proof of proper base change at the same dimension.
  \item Localization and Kunneth as derived consequences, not new axioms.
  \item Coherence tests for identity maps, composites, open immersions, proper
  maps, and cartesian rectangles.
\end{enumerate}
The formal statement of a comparison theorem should always contain the
comparison morphism itself.  Merely asserting that two objects are
noncanonically isomorphic is insufficient for every later use in trace and
duality.
